% Corrected French for printed p.160.
% Corrections: F for two stray A's; là; rationnelles; agreement and hyphen.
\SourcePage{160}

\LayerTarget{tohoku:remark:fine-sheaves-are-soft:p2}
\noindent des faisceaux fins en termes de faisceaux mous.

\LayerTarget{tohoku:construction:godement-resolution}
\textit{Résolution de l'identité.} Pour tout faisceau $F$, soit $C^0(F)$ le
faisceau produit défini par la famille des ensembles $F(x)$, on a un
homomorphisme fonctoriel $F\longrightarrow C^0(F)$, d'ailleurs injectif. Si
on se restreint à prendre $F$ dans la catégorie des faisceaux abéliens sur
$X$, la méthode de 2.5, exemple a), permet de construire une résolution de
l'identité, $C(F)$ qui se réduit à $C^0(F)$ en dimension 0 et est défini par
\[
C^n(F)=C^0\bigl(C^{n-1}(F)/\operatorname{Im}C^{n-2}(F)\bigr)
\]
en dimension $n\geq2$. Les $C^n(F)$ sont flasques, donc
$\Gamma_\Phi$-acycliques quel que soit l'antifiltre $\Phi$ de parties
fermées, donc $\Gamma_\Phi C(F)$ est un foncteur résolvant pour
$\Gamma_\Phi$, et on a par suite
\[
H^n_\Phi(X,F)=H^n\bigl(\Gamma_\Phi C(F)\bigr).
\]
$C(F)$ est appelée la \textit{résolution canonique de $F$}, (introduite et
utilisée systématiquement par Godement).

\LayerTarget{tohoku:construction:cartan-resolution}
Si la famille $\Phi$ est paracompactifiante, on trouve un autre foncteur
résolvant pour $\Gamma_\Phi$ en prenant une résolution fixe $C$ du faisceau
constant $\mathbf Z$ par des faisceaux \textit{fins et sans torsion}, et
prenant pour tout $F$ le complexe de faisceaux $F\otimes C$. C'est là une
résolution de $F$ (parce que $C$ est sans torsion), c'est un foncteur exact en
$F$ (même raison), de plus les $F\otimes C^n$ sont aussi fins donc
$\Gamma_\Phi$-acycliques ; donc le foncteur $\Gamma_\Phi(F\otimes C)$ est un
foncteur résolvant pour $\Gamma_\Phi$, et on a par suite
\[
H^n_\Phi(X,F)=H^n\bigl(\Gamma_\Phi(F\otimes C)\bigr).
\]
$F\otimes C$ sera appelée la résolution de Cartan de $F$. On se rappellera
qu'elle n'est utilisable que si $\Phi$ est paracompactifiante.

\LayerTarget{tohoku:example:irreducible-constant-sheaves}
\textit{Un exemple amusant.} Un espace $X$ est dit \textit{irréductible} s'il
n'est pas réunion de deux parties fermées distinctes de lui-même, c'est-à-dire
si l'intersection de deux parties ouvertes non vides est non vide ; il revient
encore au même de dire que toute partie ouverte de $X$ est connexe. Alors tout
faisceau \textit{constant} $F$ sur $X$ est évidemment flasque (la réciproque
étant d'ailleurs vraie si $X$ est connexe, comme on voit en prenant un
faisceau constant à fibre non réduite à un point). \textit{En particulier, si
$F$ est un faisceau constant de groupes abéliens sur l'espace irréductible
$X$, on a $H^p(X,F)=0$ pour $p>0$.}

\LayerTarget{tohoku:section:3.4}
\textbf{3.4. Applications à des questions de relèvement de groupe
structural\textsuperscript{8)}.} Soit $X$ une variété algébrique irréductible
sur un corps $k$ algébriquement clos (voir [15, Chap. II], dont nous suivons
la terminologie), $\mathbf O$ son faisceau d'anneaux locaux (=faisceau des
germes de fonctions régulières sur $X$), $\mathbf K$ le faisceau des germes de
fonctions rationnelles sur $X$. $\mathbf K$ est un faisceau constant
(loc. cit. prop. 9). Soient $\mathbf O^*$ et $\mathbf K^*$ resp. les
sous-faisceaux de $\mathbf O$ et $\mathbf K$ formés des germes inversibles ;
évidemment, $\mathbf K^*$ est encore un faisceau constant de groupes
abéliens, et $\mathbf O^*$ en est un sous-faisceau.

\LayerTarget{tohoku:construction:divisor-sequence:p1}
Le faisceau quotient $\mathbf K^*/\mathbf O^*=\mathbf D$ est le faisceau des
\textit{germes de diviseurs localement principaux sur $X$}, et coïncide avec
le faisceau des \textit{germes de diviseurs sur $X$} si les anneaux locaux
$\mathbf O(x)$ de $X$ sont factoriels (par exemple si $X$ est sans
singularités), ce que nous supposerons désormais. D'autre part, il est
immédiat que le faisceau $\mathbf D$ des germes de diviseurs sur

\LayerTarget{tohoku:footnote:8}
\begin{quote}
\footnotesize
8) La lecture du présent numéro est inutile pour la compréhension de la
suite.
\end{quote}

\TohokuContinuation{161}

